%% kapa: exam=1 round=19 year=2008 period=2 form=A subject=economics subject_ko=경제원론 seq=26 number=26
%% source: Q-Net 감정평가사 기출문제 (cst00309) artlSeq=1001962 file=extracted/1008/1001962/2008년 제19회 감정평가사 1차시험 기출문제/2교시.hwp
%% conversion: kordoc 4.24.51 -> tools/gen1.py (auto); review=proofread-sonnet; flags=none
\begin{question}{26}
경제정책에 대한 학파 간 차이를 옳게 설명한 것을 모두 고르면?
\begin{framedtext}
ㄱ. 새고전학파(new classical)에서는 예측되지 못한 정책만이 효과가 있다고 하였다. \par
ㄴ. 케인즈학파는 경기변동의 주요 원인을 총수요의 변동보다는 총공급의 변동으로 설명하였다. \par
ㄷ. 새케인즈학파(new Keynesian)는 메뉴비용이나 명목임금계약 등의 존재로 인하여 총수요 관리정책이 효과가 있다고 주장하였다. \par
ㄹ. 통화론자들은 화폐수요함수가 불안정하여 준칙에 따른 통화정책을 수행해야 한다고 주장하였다. \par
\end{framedtext}
\begin{choices}
  \item ㄱ, ㄴ
  \item ㄱ, ㄷ
  \item ㄴ, ㄹ
  \item ㄷ, ㄹ
  \item ㄱ, ㄷ, ㄹ
\end{choices}
\end{question}
\answer{2}
